\documentclass[11pt]{article}
\usepackage[margin=0.95in]{geometry}
\usepackage{xcolor}
\usepackage{booktabs}
\usepackage{enumitem}
\usepackage{titlesec}
\usepackage{microtype}
\usepackage[colorlinks=true,urlcolor=accent,linkcolor=accent]{hyperref}
\usepackage{tcolorbox}
\usepackage{needspace}

\definecolor{accent}{RGB}{72,50,160}
\definecolor{done}{RGB}{25,110,88}
\definecolor{act}{RGB}{160,95,10}
\definecolor{inkmute}{RGB}{100,96,115}
\definecolor{rulegrey}{RGB}{215,212,225}
\definecolor{boxbg}{RGB}{248,246,252}

\setlength{\parindent}{0pt}
\setlength{\parskip}{0.55em}
\setlength{\emergencystretch}{3em}

\titleformat{\section}{\large\bfseries\color{accent}}{}{0pt}{}
\titlespacing*{\section}{0pt}{1.5em}{0.6em}

% item heading: \itemhead{ID}{status}{statuscolor}{title}
\newcommand{\itemhead}[4]{%
  \needspace{4\baselineskip}%
  \vspace{0.5em}%
  {\bfseries\color{accent}#1}\hfill{\small\bfseries\color{#3}\MakeUppercase{#2}}\\[-0.4em]
  \textcolor{rulegrey}{\rule{\linewidth}{0.4pt}}\\[-0.2em]
  {\bfseries #4}\par\vspace{-0.15em}%
}

\newcommand{\quo}[1]{{\small\itshape\color{inkmute}``#1''}\par}
\newcommand{\where}[1]{%
  \vspace{0.15em}%
  {\small\textcolor{inkmute}{\textbf{Where:}} #1}\par
}

\title{\vspace{-2.2em}\bfseries Decision-Letter Compliance Map\\[0.35em]
\large\mdseries ci-2026-009433.R1 \quad$\cdot$\quad ``When Do Simple Models Win?''}
\author{}
\date{}

\begin{document}
\maketitle
\vspace{-3.2em}

\textcolor{rulegrey}{\rule{\linewidth}{1pt}}

{\small Every point in Prof.\ Kitchin's decision letter, mapped to the exact file, section and page
where it was addressed. Page numbers refer to the PDFs in \texttt{jcim\_secondrevision\_submission/}.
Audited 27 July 2026. \textbf{15 of 19 items complete and verified; 4 need an author action.}}

{\footnotesize\color{inkmute} Internal working reference --- not a submission file. Do not upload.}

\begin{tcolorbox}[colback=boxbg,colframe=act,boxrule=1pt,left=8pt,right=8pt,top=6pt,bottom=6pt,
                  title={\bfseries Four items only you can complete}, coltitle=white, colbacktitle=act]
Everything in the manuscript files is done. These four live in the submission form or need your
factual decision --- none can be settled from the \LaTeX{} sources.
\begin{enumerate}[leftmargin=1.4em,itemsep=3pt,topsep=4pt]
  \item \textbf{ORCID iDs} --- each author validates their own iD on their ACS user profile before
        the revision can be submitted. Nothing goes in the manuscript file.
  \item \textbf{Funding in the submission form} --- the manuscript already declares no specific grant
        funding (p.~48). ACS requires the same declaration in the form as well as the file.
  \item \textbf{English-language editing} --- the one item with \emph{no response anywhere} in the
        package. Needs your call on what is factually true (see p.~3).
  \item \textbf{Cover art} --- optional and opt-in. Nothing prepared; declining costs nothing.
\end{enumerate}
\end{tcolorbox}

\section{Editor's four technical issues}
{\small\color{inkmute} Answered in Part A of the response letter, pages 2--3.}

\itemhead{E1}{done}{done}{Table 5 caption recovery percentage}
\quo{The numbers in Table 5 caption don't appear to be consistent (the recovery \%) with the manuscript.}
The caption still carried the \textbf{93.6\% recovery} figure from the original submission --- the very
number retracted in the first-round response to R1-15. The Mamede photosafety set is read directly from
the published Supporting Information, not rebuilt by a Reaxys re-query, and RDKit canonicalises 74{,}783
of 74{,}784 entries ($\sim$100\%). The main text was fixed in round~1; the caption was missed. It now reads
``a reconstruction from Mamede's published Supporting Information (74{,}783 of 74{,}784 compounds
canonicalised).''
\where{Clean manuscript \textbf{p.~35} (Table 5 caption) $\cdot$ Marked \textbf{p.~40} $\cdot$
Response letter \textbf{p.~2} $\cdot$ SI \textbf{p.~33} (Section S17, full provenance)}

\itemhead{E2}{done}{done}{Regenerate Figures 5, 6 and 9}
\quo{it appears that figures 5, 6 and 9 need to be regenerated with the updated results}
Correct --- all three had been carried over from the v2 submission and never regenerated after the v3
(Greenman--Song) retrain. All three now come from the same v3 cross-validated predictions that underlie
Table~2. One subtlety is now stated explicitly so it is not mistaken for a fresh discrepancy:
\textbf{Figures 5 and 6 pool test predictions across the five folds} and report a single pooled RMSE,
whereas \textbf{Table 2 and Figure 9 report the mean across folds}. The two summaries agree to within
0.2~nm for every model, and the Figure~5 caption says which it is.
\where{Clean manuscript \textbf{p.~28} (Figs 5, 6), \textbf{p.~31} (Fig 9) $\cdot$ Marked
\textbf{pp.~32, 33, 36} $\cdot$ Response letter \textbf{p.~2} $\cdot$ Scripts
\texttt{make\_parity\_plot.py}, \texttt{regenerate\_figures.py}}

\itemhead{E3}{done}{done}{Abstract contradicts Methods on hyperparameter tuning}
\quo{The abstract suggests no hyperparameter tuning was used, but the methods suggest there was via a grid search}
Genuinely contradictory, and the abstract was the inaccurate one --- RF's tree count and per-split feature
fraction were selected by a 432-configuration grid search. Reworded in three places to describe the tuning
accurately while keeping the intended compute-cost contrast: RF needs a lightweight one-time CPU grid
search, not GPU-based optimisation.
\where{Clean manuscript \textbf{p.~1} (abstract, now ``with only lightweight grid-search tuning (no GPU)''),
\textbf{p.~4} (Introduction contributions bullet), and the Model Interpretability and Practical
Considerations section $\cdot$ Marked \textbf{p.~4} $\cdot$ Response letter \textbf{p.~3}}

\itemhead{E4}{done}{done}{Figure 10 consistency with Table 6}
\quo{Review if Figure 10 is consistent with Table 6.}
A real inconsistency. Figure~10 computed saliency and its per-molecule prediction subtitles from the
\textbf{tuned} BiGRU ensemble while Table~6 reports the \textbf{default} ensemble, so the two disagreed ---
the old caption showed Octocrylene at $+61$~nm against Table~6's $+43$~nm. Figure~10 was regenerated from
the same default-BiGRU 5-fold ensemble Table~6 uses. Per-molecule predictions now match exactly (Sinapic
Acid 318, Amiloxate 318, Padimate~O 330, Octisalate 333, Sulisobenzone 346, Cinnamic Acid 353~nm) and the
ensemble MAE of 28.6~nm matches Table~7. Caption thresholds updated
($<$22 $\rightarrow$ $<$20~nm well-predicted; $>$28 $\rightarrow$ $>$25~nm poorly predicted).
\where{Clean manuscript \textbf{p.~34} (Fig 10) against \textbf{p.~38} (Tables 6, 7) $\cdot$
Marked \textbf{p.~39} against \textbf{p.~42} $\cdot$ Response letter \textbf{p.~3}}

\section{Submission requirements}

\itemhead{TOC graphic}{done}{done}{Table of Contents graphic}
\quo{NEED Table of Contents (TOC) Graphic \ldots{} should appear on the last page of your manuscript and should not have a caption.}
Present and compliant on all three counts. It sits on the \textbf{last page} of the manuscript, carries
\textbf{no caption} (achemso's \texttt{tocentry} environment prints only the standard ``TOC Graphic''
heading), and the artwork is 975\,$\times$\,525~px --- exactly \textbf{3.25\,$\times$\,1.75 inches at
300~dpi}, the ACS minimum size at full resolution. Verified as a genuinely embedded 300~ppi RGB image,
not a missing-file placeholder.
\where{Clean manuscript \textbf{p.~56} (the final page) $\cdot$ Source
\texttt{benchmark\_paper\_jcim\_revised.tex} lines 59--61 $\cdot$ Artwork \texttt{results/toc\_graphic.png}}

\itemhead{Files a--d}{done}{done}{The four required resubmission files}
\quo{a) A point-by-point response \ldots{} b) A clean, unmarked copy \ldots{} c) A marked copy \ldots{} Upload this file as Supporting Information for Review. d) A clean, unmarked Supporting Information for Publication file.}
All four are built and verified, plus the cover letter. Every document compiles with \textbf{zero undefined
references or citations}. The clean manuscript, marked manuscript and Supporting Information reproduce
\textbf{byte-for-byte} from their committed sources, which proves none of them is a stale build. Full
manifest on the last page.

\itemhead{Funding}{your action}{act}{Funding sources}
\quo{Confirm all sources of funding for ALL authors relevant to this manuscript are included in BOTH the submission form and in the manuscript file.}
The \textbf{manuscript half is done}: the Acknowledgement declares that the research received no specific
grant from funding agencies in the public, commercial or not-for-profit sectors. A nil declaration
satisfies the requirement. \textbf{What is left is the form half} --- ACS asks for the same declaration
inside Paragon Plus, and that field cannot be set from the \LaTeX{} files.
\where{Clean manuscript \textbf{p.~48} (Acknowledgement) $\cdot$ Marked \textbf{p.~54} $\cdot$
Cover letter \textbf{p.~2}}

\itemhead{ORCID}{your action}{act}{Validated ORCID iDs}
\quo{Authors submitting manuscript revisions are required to provide their own validated ORCID iDs before completing the submission.}
Purely a submission-system task --- ORCID iDs attach to ACS user profiles, never to the manuscript file, so
there is correctly nothing to change in the \LaTeX{}. Each author validates their own iD; you cannot supply
one on a co-author's behalf. The cover letter already tells the editor these are provided through the
system, so make sure that is true before you submit.
\where{ACS Paragon Plus $\rightarrow$ Email/Name section of each author's account $\cdot$
Cover letter \textbf{p.~2}}

\itemhead{English}{needs your call}{act}{Quality of written English}
\quo{The reviewers indicate that the quality of the written English could be improved. Please consider using a native English speaker or an English-language editing service to assist with this.}
\textbf{This is the one item with no response anywhere in the package} --- not in the cover letter, not in
the response document. It appears in the standard block of the letter rather than in either reviewer's
numbered comments, and neither reviewer raised language in their own text, so it is plausibly boilerplate.
Even so, it was explicitly asked, and a single sentence in the cover letter closes it.

I have deliberately not written that sentence, because it is a factual claim about what you did and only
you can make it truthfully. Two honest options: state that the manuscript was reread and copy-edited for
clarity and grammar during this revision, which the prose-cleanup work supports; or, if you used a service
or a native-speaker colleague, name that. Do not claim a professional service that was not used.
\where{Would go in \texttt{cover\_letter\_revised.tex}, after the Reviewer~1 paragraph, before the file list}

\itemhead{Cover art}{optional}{act}{Cover art}
\quo{JCIM is now accepting images to appear on the front cover \ldots{} please upload it and a brief (80 word) description.}
An open invitation, not a requirement --- it uses ``if you have,'' and declining costs nothing. Nothing has
been prepared and no cover art is referenced anywhere in the package. If you did want to pursue it, the TOC
graphic is the natural starting point, but it would need re-rendering to the cover specification and an
80-word caption, uploaded under the \emph{Cover Art} and \emph{Cover Art Caption} designations.

\itemhead{iThenticate}{no action}{inkmute}{Similarity screening}
Informational --- ACS runs this automatically and nothing is required of you. Worth knowing that a revision
legitimately overlaps heavily with its own earlier version, which the screening accounts for.

\section{Reviewer 1's seven suggestions}
{\small\color{inkmute} Answered in Part A of the response letter, pages 4--7. All seven complete.}

\itemhead{R1 $\cdot$ 1}{done}{done}{Reviewer comments were paraphrased, not verbatim}
\quo{This is true for my comments R1-1 to R1-11, but R1-12 to R1-19 are paraphrased \ldots{} The authors should double check that the paraphrased reviewer comments are the only things that the AI didn't copy correctly.}
Every reviewer-comment block in Part~B was replaced with the exact text of the original decision letter.
The requested double-check turned up \textbf{more than the reviewer had spotted}: Reviewer~2's comments
R2-1, R2-2, R2-4, R2-5, R2-6 and R2-7 had also been paraphrased. Those were restored too, and the response
says so plainly rather than quietly fixing them.

On the wider audit, the response states that the only substantive inaccuracies found are the two already
reported in the first-round response --- the ``variance is intrinsic to Chemprop'' claim, corrected by the
R1-12 HPO probe, and the 93.6\% Reaxys recovery claim, corrected in R1-15. The inference about AI
assistance is confirmed directly rather than deflected, and cross-referenced to the disclosure in point~5.

\textit{Note:} the response document now covers the second-round comments only --- the first-round
point-by-point (former Part~B) was removed to keep it focused. The corrected verbatim blocks therefore
are not reprinted; the response says so and offers them on request.
\where{Response letter \textbf{p.~4} (Comment 1)}

\itemhead{R1 $\cdot$ 2}{done}{done}{Supporting Information sections were not numbered}
\quo{References to the supporting information from the main text are of the form ``Supporting Information Section S7'', but the sections in the actual supporting information file are not numbered.}
Section numbering is enabled and prefixed, so the SI now carries \textbf{21 numbered sections, S1 through
S21}, matching how the main text refers to them. All 28 \texttt{Section S$n$} references in the manuscript
were independently checked against the SI's actual section list --- every one resolves to the topically
correct section, with no off-by-one drift.
\where{Source \texttt{supporting\_information\_revised.tex} line 44:
\texttt{\textbackslash renewcommand\{\textbackslash thesection\}\{S\textbackslash arabic\{section\}\}}
$\cdot$ Response letter \textbf{p.~4}}

\itemhead{R1 $\cdot$ 3}{done}{done}{v2 and v3 dataset versions need clearer definition}
\quo{The new ``v2'' and ``v3'' descriptions of datasets used through the main text and SI could be more clearly defined, especially in their first usage. I also encourage the authors to consider using more descriptive titles.}
Both halves are addressed. Each version is now defined at first use and given a descriptive name rather
than a bare version number --- v2 as the 18{,}755-row first-occurrence-dedup set, v3 as the 18{,}415-row
\textbf{Greenman--Song} duplicate-handling set. The abstract itself now names the protocol instead of
saying ``v3.'' A dedicated SI section compares the two side by side.
\where{Clean manuscript \textbf{p.~1} (abstract names the protocol), \textbf{p.~21} (Table 2 dagger
footnote flags the two v2 rows) $\cdot$ SI \textbf{p.~16} (Section S10) $\cdot$ Response letter \textbf{p.~5}}

\itemhead{R1 $\cdot$ 4}{done}{done}{Inter-revision phrasing in captions}
\quo{Some places in the manuscript mention how things were changed from one revision to the next, which is unconventional \ldots{} Table S24 caption also appears to be missing a word: ``(adopted in this revision per )''}
Both flagged captions are fixed, and the truncated one had an underlying cause worth noting: the missing
word was a \textbf{broken citation} that rendered as an empty reference, not a typo. That citation is
repaired, so the caption is complete rather than merely reworded. Both the manuscript and the SI were swept
for any remaining revision-relative phrasing --- none survives.
\where{SI Table S16 and Table S24 captions $\cdot$ Response letter \textbf{p.~5} $\cdot$ Verified: zero
occurrences of ``in this revision'' / ``the revised manuscript'' remain in either file}

\itemhead{R1 $\cdot$ 5}{done}{done}{AI-assistance disclosure understates the scope}
\quo{If AI assisted with any of the manuscript writing/revising in addition to formatting, I encourage the authors to add this to the disclosure \ldots{} the CLAUDE.md file seems to indicate that AI was used more heavily than just for code debugging/review.}
The disclosure was broadened rather than defended. It now reads that Claude (Anthropic) assisted with
\textbf{code development and debugging, data-analysis scripting, and editing of the manuscript} ---
explicitly covering the manuscript writing the reviewer inferred from the commit history and CLAUDE.md.
The authorship boundary is stated in the same sentence: all scientific content, analysis, experimental work
and interpretation are the authors' own, and the authors take responsibility for the final manuscript.
AI assistance with the response document itself is disclosed in the response letter (Comments~1 and~5)
rather than in the published Acknowledgement.
\where{Clean manuscript \textbf{p.~48} (Acknowledgement) $\cdot$ Marked \textbf{p.~54} $\cdot$
Response letter \textbf{p.~6}}

\itemhead{R1 $\cdot$ 6}{done}{done}{Repository organisation and the missing LICENSE}
\quo{I suggest that the authors clean up the repository a bit by creating more subdirectories \ldots{} while the authors specify in a couple places in the README that their project is published under the MIT License, there is no LICENSE file present.}
A dedicated, cleanly organised repository was created for the manuscript: \textbf{ML\_UV\_models\_jcim}.
Verified live --- it is public, the root holds just four directories (\texttt{data}, \texttt{results},
\texttt{scripts}, \texttt{uvml}) plus README, pyproject, requirements and gitignore, and a \textbf{LICENSE
file is present and detected by GitHub as the MIT License}. The scattered \texttt{.tex} files and images the
reviewer objected to are no longer in the root. The original development repository stays public so the full
commit history the reviewer inspected remains available; the manuscript cites both and explains why each
exists.
\where{Clean manuscript \textbf{p.~48} (Data and Software Availability) $\cdot$ Response letter \textbf{p.~6}
$\cdot$ Verified live: \texttt{ML\_UV\_models\_jcim} (200), \texttt{ML\_UV\_models} (200),
Zenodo \texttt{10.5281/zenodo.20600225} (200)}

\itemhead{R1 $\cdot$ 7}{done}{done}{Figure 2(b) implies sequential D-MPNN encoders}
\quo{The Chemprop D-MPNN part of this figure is misleading because it implies the Solute and Solvent D-MPNNs are sequential rather than parallel. This is explained correctly in the caption, but I think the figure should be updated accordingly.}
The panel was redrawn so the solute and solvent encoders appear as genuine parallel branches converging on a
concatenation, matching what the caption always said. The caption was tightened to match the new drawing ---
``two parallel directed message-passing encoders \ldots{} concatenates their mean-pooled representations
before a feed-forward regression head'' --- and the closing sentence now attributes solvent handling in
panel~(b) to the parallel-encoder multicomponent configuration.
\where{Clean manuscript \textbf{p.~14} (Figure 2, panel b) $\cdot$ Marked \textbf{p.~15} $\cdot$
Response letter \textbf{p.~7}}

\section{Caught during the pre-submission audit}
{\small\color{inkmute} Not requested by anyone --- found while verifying the package was ready to upload.}

\itemhead{Fix}{fixed}{done}{The response letter had a broken bibliography}
The staged response-to-reviewers PDF had been compiled without a complete BibTeX pass.
\textbf{Fourteen citations rendered as literal \texttt{[?]}} --- among them the Greenman, Song and
Chemprop~v2 references that several answers depend on --- and \textbf{the References page was missing
entirely}. Given that Reviewer~1's first point was precisely about citation and quoting rigour in this
document, sending it that way would have read badly.

Rebuilt from source: all fourteen citations resolve to numbered entries and the References list is restored.
Old and new were diffed to confirm nothing else changed. The cover letter was rebuilt at the same time so
both now carry the 27 July submission date rather than 15 July. The document was subsequently trimmed to
the second-round responses alone (the first-round point-by-point was removed), taking it from 46 pages to~8.

\textit{Worth carrying forward:} a grep for \texttt{??} finds undefined \emph{references} but not undefined
\emph{citations}, which render as \texttt{[?]}. The first scan came back clean and missed this; it only
surfaced on recompiling from source and diffing against the staged file.

\section{Reading the marked copy}
The marked manuscript carries two rounds of changes in two separate colour pairs, so the editor can isolate
what is new since the last decision. \textbf{Round-2 edits --- everything in this document --- are green and
orange.}

\vspace{-0.3em}
\begin{center}
\begin{tabular}{@{}llr@{}}
\toprule
\textbf{Colour} & \textbf{Meaning} & \textbf{Count} \\
\midrule
Blue & Round-1 additions & 126 \\
Red & Round-1 deletions & 85 \\
\textbf{Green} & \textbf{Round-2 additions} & \textbf{13} \\
\textbf{Orange} & \textbf{Round-2 deletions} & \textbf{10} \\
\bottomrule
\end{tabular}
\end{center}

Marking is colour-only rather than struck through, because soul's \texttt{\textbackslash st} cannot span
paragraph breaks or table cells and would have broken several of the edited captions.

\section{Upload manifest}
{\small\color{inkmute} Everything in \texttt{jcim\_secondrevision\_submission/}, with its ACS file designation.}

\vspace{-0.3em}
\begin{center}
{\small
\begin{tabular}{@{}p{6.3cm}p{4.6cm}r l@{}}
\toprule
\textbf{File} & \textbf{ACS designation} & \textbf{pp.} & \textbf{State} \\
\midrule
\texttt{\ldots\_revised.pdf} & Manuscript & 56 & Verified \\
\texttt{\ldots\_revised\_marked.pdf} & Supp.\ Info.\ for Review Only & 62 & Verified \\
\texttt{\ldots\_response\_to\_reviewers.pdf} & Response to Reviewers & 8 & \textbf{Rebuilt} \\
\texttt{\ldots\_supporting\_information.pdf} & Supp.\ Info.\ for Publication & 47 & Verified \\
\texttt{\ldots\_cover\_letter.pdf} & Cover Letter & 2 & \textbf{Rebuilt} \\
\bottomrule
\end{tabular}}
\end{center}

{\footnotesize\color{inkmute}
All files carry the \texttt{When\_do\_simple\_models\_win\_} prefix. Clean manuscript, marked manuscript and
Supporting Information reproduce byte-for-byte from their committed \LaTeX{} sources. Every external link
was checked live and resolves.}

\end{document}
